
\srcspaced{下面还要证明：对无穷多个 $n$，存在次数为 $2^n$ 的循环域，它们是四元数群的极小分裂域。}

设 $p$ 是一个有理素数，使得对某个适当的 $r>0$，有
\[
        (*)\qquad 2^r+1\equiv0\pmod p
\]
成立；并设 $2^n$ 是整除 $p-1$ 的最高的 $2$ 的幂。

我们先证明，对无穷多个 $n$ 都存在素数 $p$。为此，设 $k$ 是任意一个正有理整数，$p$ 是 $2^{2^k}+1$ 的一个素因子。于是，关于 $p$ 所要求的同余在 $r=2^k$ 时成立。此外，$2\pmod p$ 的指数为 $2^{k+1}$，因而对整除 $p-1$ 的最高的 $2$ 的幂 $2^n$，数 $n$ 大于 $k$。所以无论如何，都有任意大的 $n$ 与某些素数 $p$ 相对应。

设 $\varepsilon$ 是一个本原 $p$ 次单位根；我们现在证明，在 $P(\varepsilon)$ 中，$-1$ 可以表示为两个平方之和
\[
        -1=a^2+b^2=(a+ib)(a-ib)\qquad (a,b\text{ 是 }P(\varepsilon)\text{ 中的数})
\]
。这等价于：在 $(4p)$ 次单位根的域 $P(\varepsilon,i)=P(\varepsilon\cdot i)$ 中，存在关于 $P(\varepsilon)$ 的相对范数恰为 $-1$ 的数。

数 $1+i\varepsilon^\nu$（$\nu=1,2,\ldots,p-1$）的相对范数为 $(1+i\varepsilon^\nu)(1-i\varepsilon^\nu)=1+\varepsilon^{2\nu}$；因此，所有数 $1+\varepsilon^\nu$（$\nu=1,2,\ldots,p-1$）都是 $P(\varepsilon i)$ 中的数关于基域 $P(\varepsilon)$ 的相对范数。所以同样的结论也适用于
\[
\Pi=\prod_{\nu=0}^{r-1}(1+\varepsilon^{2^\nu})
=(1+\varepsilon)(1+\varepsilon^2)(1+\varepsilon^4)\cdots(1+\varepsilon^{2^{r-1}})
=\sum_{\lambda=0}^{2^r-1}\varepsilon^\lambda.
\]
若再向 $\Pi$ 中添入 $\varepsilon^{2^r}=\varepsilon^{-1}=\varepsilon^{p-1}$（由 $(*)$），则 $\Pi+\varepsilon^{p-1}$ 成为若干个整段 $1+\varepsilon+\varepsilon^2+\cdots+\varepsilon^{p-1}=0$ 之和，因而
\[
        \Pi=-\varepsilon^{-1}.
\]
现在，$\varepsilon$ 也是 $P(\varepsilon i)$ 中一个数的相对范数，即 $\varepsilon^{(p+1)/2}$ 的相对范数。因此，类似的结论适用于
\[
        \Pi\cdot\varepsilon=-1.
\]
所以 $-1$ 可以在 $P(\varepsilon)$ 中表示为两个平方之和；\footnote{对于形如 $8n\pm3$ 的素数，这一论证已见于 E. Landau, Über die Darstellung definiter Funktionen durch Quadrate, Math. Ann. Bd. 62. S. 272--285。} $P(\varepsilon)$ 是四元数群的一个分裂域。

设 $K$ 是 $P(\varepsilon)$ 中所含的次数为 $2^n$ 的循环域；我们断言，$K$ 也是一个分裂域。\footnote{把次数为 $2^n$ 的循环分裂域选作分圆域的子域，是仿照 Hasse 的例子。} 设 $m$ 是四元数除环关于以 $K$ 为基域时的指数，则 $m=1$ 或 $m=2$。但是，因为存在相对于 $K$ 的相对次数为奇数 $(p-1)/2^n$ 的分裂域 $P(\varepsilon)$，第二种情形不可能；所以 $m=1$。而这正是说，$K$ 是一个分裂域。

因此，对无穷多个 $n$，存在次数为 $2^n$ 的循环域，它们是极小分裂域。

\begin{center}
\rule{7em}{0.4pt}
\end{center}

\clearpage
\setcounter{footnote}{0}

